\documentclass[12pt]{article}
\usepackage[T1]{fontenc}
\usepackage[utf8]{inputenc}
\usepackage{lmodern}
\usepackage{textcomp}
\usepackage{enumitem}
\usepackage{geometry}
\geometry{margin=1in}
\begin{document}
\begin{center}
Answer Key - Health Sciences - Undergraduate Level Assessment 14 \
\small (Answers are ordered exactly as in the question paper.)
\end{center}

\section*{Multiple Choice}
\begin{enumerate}[label=\arabic*.]
\item \textbf{Answer:} A
\\ \textit{Options:} A) A cytokine storm; B) Antibody mediated reactions; C) A response of killer T cells (CD-8); D) A decline in function of the immune system
\\ \textit{Note:} A cytokine storm occurs when the immune system overreacts to a viral infection, leading to an excessive release of pro-inflammatory cytokines that can result in tissue damage and severe immune pathology.
\item \textbf{Answer:} D
\\ \textit{Options:} A) Dendrite; B) Glial cell; C) Nerve center; D) Synapse
\\ \textit{Note:} The synapse is the specialized junction between two neurons where the transmission of an impulse occurs through the release of neurotransmitters, facilitating communication within the nervous system.
\item \textbf{Answer:} A
\\ \textit{Options:} A) Collagen; B) Heparin; C) Lipocyte; D) Melanin
\\ \textit{Note:} Collagen is the primary structural protein in the dermis, providing strength and elasticity to the skin, which classifies it as a fibrous protein.
\item \textbf{Answer:} A
\\ \textit{Options:} A) Because they have no nucleic acid; B) They require a helper virus; C) Only replicate in dividing cells; D) Can integrate into host chromosomes
\\ \textit{Note:} The statement that parvoviruses have no nucleic acid is incorrect; rather, parvoviruses possess single-stranded DNA, which is essential for their replication and pathogenicity, making them impactful parasites due to their ability to infect and exploit host cells effectively.
\item \textbf{Answer:} D
\\ \textit{Options:} A) Food and Drug Administration reforms; B) Easier access to investigational drugs; C) Changes in the way medicine was practiced in the U.S.; D) All of the above
\\ \textit{Note:} AIDS activism in the U.S. resulted in increased awareness of the disease, significant changes in public health policy, and the mobilization of communities to advocate for research funding and treatment access, making "All of the above" the correct choice as it encompasses these multifaceted impacts.
\end{enumerate}

\section*{True / False}
\begin{enumerate}[label=\arabic*.]
\setcounter{enumi}{5}
\item \textbf{Answer:} False
\\ \textit{Options:} A) True; B) False
\\ \textit{Note:} The muscles of the soft palate are primarily innervated by the vagus nerve, specifically through the pharyngeal plexus, rather than the facial or glossopharyngeal nerves.
\item \textbf{Answer:} True
\\ \textit{Options:} A) True; B) False
\\ \textit{Note:} Chemotherapy can effectively manage herpes infections by alleviating symptoms and minimizing viral load, although it is primarily used for cancer treatment and antiviral medications are typically more effective for herpes.
\item \textbf{Answer:} False
\\ \textit{Options:} A) True; B) False
\\ \textit{Note:} Reliability refers to the consistency and stability of a measure over time, while validity pertains to whether the measure accurately captures the intended concept or construct.
\item \textbf{Answer:} False
\\ \textit{Options:} A) True; B) False
\\ \textit{Note:} A dished face profile is typically associated with conditions like craniosynostosis or congenital syndromes, rather than an enlarged frontal bone due to hydrocephaly, which more commonly results in a prominent forehead and increased intracranial pressure.
\item \textbf{Answer:} True
\\ \textit{Options:} A) True; B) False
\\ \textit{Note:} Research indicates that pessimism is associated with negative health outcomes and higher mortality rates, which may lead to a higher prevalence of optimism among surviving older adults.
\end{enumerate}

\section*{Long Form Response}
\begin{enumerate}[label=\arabic*.]
\setcounter{enumi}{10}
\item \textbf{Answer:} The adenovirus virion exhibits a unique icosahedral shape, which is a geometric structure composed of 20 triangular faces that provide a stable and symmetric form. This configuration is critical for protecting the viral genetic material and enables efficient assembly and stability in the extracellular environment. Additionally, adenoviruses possess slender fibers that extend from the vertices of the icosahedron, which play a crucial role in the virus's infectivity. These fibers facilitate attachment to host cell receptors, promoting viral entry and enhancing the efficiency of infection. Together, the icosahedral structure and slender fibers contribute to the adenovirus's ability to withstand environmental challenges while effectively targeting and infecting host cells.
\\ \textit{Note:} The answer correctly highlights the structural features of the adenovirus, specifically its icosahedral shape for stability and protection of genetic material, as well as the slender fibers that enhance infectivity by facilitating attachment to host cells.
\item \textbf{Answer:} The spread and infection of viruses are influenced by several key factors, with travel being one of the most significant. Travel facilitates the rapid movement of individuals across geographical boundaries, allowing viruses to hitch rides to new locations and populations. This is particularly evident in the global interconnectedness fostered by air travel, where infected individuals can transmit viruses to new areas before symptoms even manifest. Furthermore, urbanization and the density of populations in transit hubs can exacerbate the likelihood of outbreaks. Overall, understanding the role of travel in virus transmission is crucial for public health strategies aimed at controlling and preventing infectious diseases.
\\ \textit{Note:} The answer correctly identifies travel as a crucial factor in the spread of viruses by illustrating how it enables the swift transmission of infections across borders, particularly through air travel, which can lead to outbreaks in new locations before symptoms appear.
\end{enumerate}

\vfill
\noindent\textit{End of Answer Key}
\end{document}